\section{Finite Fourier coefficients in a logarithmic variable}
\label{sec:fourier-coefficients}

The polynomial coefficient algebra extends naturally to finite sums
\[
 B(L)=\sum_{\omega\in\Omega}P_\omega(L)e^{i\omega L},
 \qquad L=\log u,\qquad \Omega\subset\R\text{ finite},
\]
where every $P_\omega$ is a polynomial and the frequencies are exact real
numbers. This includes $\sin(\log u)$, several incommensurable frequencies,
and polynomial amplitudes multiplying those oscillations. The source powers
and the Fourier frequencies have different roles: source weights determine
asymptotic truncation, while frequencies label coefficient modes. No Fourier
mode is discarded merely because its frequency is large.

\subsection{Coefficient arithmetic and reality}

Multiplication convolves frequencies, and the Euler derivative is
\[
 \frac{\dd}{\dd L}\bigl(P_\omega(L)e^{i\omega L}\bigr)
 =\bigl(P_\omega'(L)+i\omega P_\omega(L)\bigr)e^{i\omega L}.
\]
Thus this algebra is closed under precisely the operations used by
Theorem~\ref{thm:coefficients}. Distinct syntactic frequencies are combined
when exact comparison proves them equal. Algebraic frequencies are reduced
before comparison. No floating-point tolerance is used to identify modes;
an undecidable comparison produces a failure.

The real-valued condition is the conjugacy relation
$P_{-\omega}(L)=\overline{P_\omega(L)}$ on real $L$. The parser checks it
under the supplied parameter assumptions. Intermediate complex coefficients
are permitted, and real trigonometric expressions are reconstructed for
display. Affine trigonometric phases, such as $\sin(2L+1)$, contribute their
constant phase factor to the polynomial amplitude.

\subsection{Inverse coefficients and independent composition}

The admitted forward family is
\[
 f(u)=y_0+a u^p\left(1+\sum_j u^{\delta_j}B_j(\log u)\right),
 \qquad ap\ne0,\quad \delta_j>0,
\]
with finitely many real Fourier-polynomial coefficients. The leading core is
a monomial with a proved nonzero real coefficient. An oscillatory leading
block, even one that happens to remain positive, requires a different inverse
normalization and is rejected explicitly by this engine.

For a multi-index $\bk$ of depth $n$ and weight
$\sigma=\bk\cdot\bdelta$, the unchanged coefficient formula is
\[
 C^{(r)}_{\bk}(L)=\frac{(-1)^n r}{p^n\bk!}
 \prod_{j=1}^{n-1}(\D+r+\sigma+pj)\prod_i B_i(L)^{k_i}.
\]
Convolution forms the product, and the Fourier Euler operation applies each
factor. The implementation preserves the individual coefficients as well as
the complete sum at a source weight.

Residual checking uses a separate composition recurrence. If
$u=z(1+U)$, then
\[
 (1+U)^\alpha B(L+\log(1+U))=\sum_{k\ge0}Q_k(L)U^k,
 \qquad Q_0=B,\qquad
 Q_{k+1}=\frac{(\alpha-k)Q_k+\D Q_k}{k+1}.
\]
This follows by Taylor expansion in the multiplicative increment; it holds
for each exponential-polynomial mode as well as for their sum. Its use in
the normalized forward equation independently checks the inverse coefficients.
For a nonunit observable power, the residual first reconstructs the source
unit by raising the observable unit to its reciprocal power.

\subsection{Envelopes, convergence, and omitted terms}

For real $L$, Fourier modes have modulus one. If
$P_\omega(L)=\sum_j c_{\omega,j}L^j$, then
\[
 |B(L)|\le\sum_{\omega,j}|c_{\omega,j}|(1+|L|)^j.
\]
This nonoscillatory envelope is stored with each retained block. Euler
differentiation preserves the polynomial degree; frequency factors affect
coefficient constants, rather than the power of the logarithmic envelope.
Consequently a multi-index coefficient has polynomial degree at most
$\sum_i k_i\deg B_i$, just as in the polynomial case.

The convergence argument also extends. After replacing each finite source
mode by an independent marker, logarithmic composition uses
\[
 e^{i\omega(L+\log W)}=e^{i\omega L}W^{i\omega},
\]
which is analytic for $W$ near $1$ with its local logarithm. The implicit
derivative at the origin is still $p\ne0$. The finite-dimensional analytic
implicit function theorem gives an absolutely convergent marker expansion
for sufficiently small markers. Along the real source slice these markers
are bounded by constants times $z^{\delta_i}(1+|\log z|)^{d_i}$ and tend to
zero. The positive source gaps also make the forward derivative asymptotic
to $ap u^{p-1}$, so the selected real inverse branch exists near the endpoint.

The complete one-step outer boundary of the retained multi-index region
therefore bounds the omitted tail. Its least weight supplies the remainder
power, and the maximum summed polynomial degree over that finite boundary
supplies a conservative logarithmic envelope. This bound need not be sharp
after cancellation, but remains valid even at zeros of every displayed
oscillatory coefficient. A declared forward remainder is transported with
the same observable-dependent exponent as Theorem~\ref{thm:transport}; its
matching derivative hypothesis remains required. Constants and source
thresholds are not computed numerical certificates.

\subsection{Interface and acceptance examples}

\begin{lstlisting}
s = AsymptoticFourierInverse[x + x^2 Sin[Log[x]],
  {x, 0}, {y, 4}];
FourierInverseResidual[s]
FourierInverseCoefficient[s, {2}]
\end{lstlisting}
The result is
\[
 y-y^2\sin L+y^3\sin L(2\sin L+\cos L)+\OO(y^4),
 \qquad L=\log y.
\]
With an additional $x^3\cos(2\log x)$, the coefficient of $y^3$ acquires
$-\cos(2L)$, combining a directly supplied frequency with frequencies
generated by convolution. The identity
\[
 \sin(\sqrt2L)\sin(\sqrt3L)
 +\tfrac12\cos\!\bigl(\sqrt{5+2\sqrt6}\,L\bigr)
 =\tfrac12\cos\!\bigl((\sqrt2-\sqrt3)L\bigr)
\]
tests exact algebraic frequency equality and cancellation.

The cutoff is exclusive in the positive target power coordinate, with the
same finite-endpoint, infinity, and observable conventions as the ordinary
engine. This first Fourier engine requires an explicit cutoff and uses the
individual Lagrange formula. Frequencies are exact numeric real constants;
arbitrary symbolic frequency relations and nonlinear logarithmic phases are
outside this interface. \wl{"MaxFrequencies"} bounds distinct modes across
the retained jet and intermediate coefficient operations, and
\wl{"MaxTerms"} bounds source indices and convolution work. Exhaustion
returns an explicit failure; there is no undocumented frequency truncation.

Regression cases compare the displayed formulas with direct trigonometric
differentiation, independently compose normalized residuals, check source and
target translations and an infinite source endpoint, and exercise observable
powers, algebraic resonance, polynomial envelopes, remainder transport, and
frequency budgets. The Fourier envelope test deliberately evaluates at a zero
of the oscillatory factor, where a sign-based or oscillatory error scale would
give a misleading result.
